\chapter{Advanced Packaging for Optical Engines}
\label{ch:packaging}
\label{sec:advanced-packaging}

This chapter is the design judgment home for short-reach optical packaging:
what 2.5D and 3D actually mean, why co-packaged optics is a placement choice
rather than a stack recipe, and how thermal partitioning, known-good die,
yield, fiber attach, and serviceability decide whether the engine ships.
Program status, shipping CPO products, and XPO hedges stay in
\Cref{app:ai-fabric-context,sec:cpo-status,sec:xpo}. External laser modules
live in \Cref{sec:elsfp}. SerDes reach classes and module styles are in
\Cref{sec:serdes-dsp,sec:conditioning,sec:pluggables}.

\textit{Read first:} 2.5D versus 3D, CPO is not automatically 3D, thermal
partitioning, and known-good-die / yield.

\textit{Deep dive:} placement ladder energy and serviceability in
\Cref{tab:placement}; ELSFP architecture in \Cref{sec:elsfp}; COUPE and
vendor programs in \Cref{sec:cpo-status}.

\keyidea{CPO means the optical engine sits in or next to the ASIC package. It
may use 2.5D, 3D, substrate-mounted engines, or an interposer. Do not equate
CPO with vertical stacking. The first packaging question is which impairment
you are buying: electrical reach, thermal access, yield, or field service.}

\section{Define the terms carefully}
\label{sec:pkg-terms}

\subsection{2.5D integration}
\label{sec:pkg-25d}

\term{2.5D} integration places multiple dies side by side
on a shared platform: a silicon interposer, an organic interposer, a bridge
die, or a redistribution layer. Communication is primarily lateral through
dense interconnect. Chiplets, HBM stacks beside logic, and photonics sitting
next to a switch ASIC are common 2.5D patterns. Electrical reaches are short
compared with a faceplate VSR run, but the dies still sit in one plane for
heat extraction and often for separate test before attach.

\subsection{3D integration}
\label{sec:pkg-3d}

\term{3D} integration stacks dies vertically using hybrid
bonding, microbumps, through-silicon vias, or face-to-face / face-to-back
bonds. Communication is vertical and very dense. Bandwidth density rises;
thermal access and test access get harder. Stack yield can multiply die
yields, so known-good-die strategy becomes critical. Optical I/O through a
tall stack is not automatic. Fiber exit, laser placement, and heat paths
must be designed into the stack, not assumed.

\subsection{Co-packaged optics}
\label{sec:pkg-cpo-def}

\term{Co-packaged optics}\firstterm{CPO} places optical engines in the same
package as the switch or accelerator ASIC, or close enough that the
electrical hop is XSR-class millimeters rather than VSR-class host traces
(\Cref{tab:placement,sec:224g}). CPO may use:

\begin{itemize}
\item 2.5D integration (engine beside the ASIC on an interposer or substrate);
\item 3D integration (electronics bonded onto photonics, or the reverse);
\item substrate-mounted optical engines without a full interposer;
\item interposer-attached photonic I/O;
\item external laser sources feeding the package (\Cref{sec:elsfp}).
\end{itemize}

The interview trap is ``CPO means 3D hybrid bonding.'' That is one recipe,
not the definition. Ask what electrical reach, thermal budget, and
serviceability the design actually bought.

\section{2.5D versus 3D}
\label{sec:pkg-25d-vs-3d}

\begin{center}
\footnotesize
\begin{tabular}{@{}p{0.22\linewidth}p{0.34\linewidth}p{0.34\linewidth}@{}}
\toprule
Topic & 2.5D & 3D \\
\midrule
Electrical reach & Short lateral links & Very short vertical links \\
Bandwidth density & High & Extremely high \\
Thermal access & Easier & Harder \\
Test access & Better & More difficult \\
Known-good-die strategy & More manageable & Critical and harder \\
Yield interaction & Dies can be screened separately & Stack yield can multiply \\
Repairability & Better & Poor \\
Optical access & Often easier & Depends strongly on stack \\
Package complexity & High & Very high \\
Best use & Chiplets, HBM, adjacent photonics & Dense logic/memory or bonded
electronics-photonics \\
\bottomrule
\end{tabular}
\end{center}

{\raggedright\footnotesize Interview memory: 2.5D is side-by-side with lateral
wiring. 3D is stacked with vertical wiring. CPO can use either. Thermal and
KGD often decide before bandwidth density does.\par}

\section{Optical-engine packaging issues}
\label{sec:pkg-engine-issues}

Once the stack style is named, the engine still has to close a list of hard
physical problems.

\begin{description}
\item[Fiber attach / FAU] Fiber-array units set coupling loss, polarization
      alignment, and lane-to-lane uniformity. A one-channel FAU miss looks
      like a dark lane after the package is sealed.
\item[Optical I/O orientation] Edge coupling, grating coupling, and vertical
      fiber exits fight different keep-out zones, lid designs, and board
      clearances.
\item[Thermal isolation] Lasers hate the ASIC hotspot. Rings and filters
      drift with local heaters and neighbor channels
      (\Cref{ch:wdm,ch:lasers}). Partition heat deliberately.
\item[Modulator temperature] Resonant modulators transfer package thermal
      gradient into wavelength and OMA control burden.
\item[PD/TIA proximity] Co-packaging the detector with the TIA cuts
      capacitance and noise (\Cref{ch:models}). Bondwire length is a noise
      budget line, not a drawing convenience.
\item[Bump pitch and RF loss] Dense electrical interfaces buy bandwidth and
      spend SI margin. Package discontinuities show up as EQ and COM
      problems (\Cref{ch:equalization}).
\item[Power delivery] Optical engines and SerDes share rails and return
      paths with a multi-hundred-watt ASIC. Supply noise becomes optical
      margin.
\item[Warpage] Large packages warp over reflow and temperature. Alignment
      and bump reliability move with the warp.
\item[Known-good die] Screen dies before you stack or co-package them when
      you can. A bad photonic engine bonded to a good switch scrapes both.
\item[Assembly yield] Optical attach, underfill, lid, and fiber steps often
      dominate module cost more than the PIC wafer.
\item[Test strategy] Define what you measure before lid attach, after fiber
      attach, and at system bring-up. Late discovery is expensive.
\item[Serviceability] Decide what fails in the field and what you can
      replace: laser module, fiber jumper, whole package, or nothing.
\end{description}

\section{Co-packaging tradeoffs}
\label{sec:pkg-cpo-tradeoffs}

Moving optics onto the package buys and spends at the same time.

\paragraph{What you buy.}
\begin{itemize}
\item Shorter SerDes reach (XSR-class) and lower electrical energy per bit
      (\Cref{tab:placement,sec:power}).
\item Higher shoreline bandwidth density than a faceplate cage farm.
\item Less host-board SI pain for the optical hop.
\end{itemize}

\paragraph{What you spend.}
\begin{itemize}
\item Thermal coupling to a large ASIC: rings walk, lasers age faster if
      left on-package, lock loops fight neighbors.
\item Package and assembly yield: one bad optical attach can scrap an
      expensive switch package.
\item Fiber management at the package edge instead of at a pluggable cage.
\item Harder test access after final assembly.
\item Field replacement: soldered engines are not hot-swappable; that is why
      many CPO recipes pull lasers into ELSFP banks (\Cref{sec:elsfp}).
\end{itemize}

XPO and pluggable hedges exist because some fleets still want faceplate
serviceability while they close CPO yield and thermal stories
(\Cref{sec:xpo,sec:pluggables}). Productization then has to prove the factory
can reproduce the package claim (\Cref{ch:productization}).

\subsection{External versus integrated lasers}
\label{sec:pkg-laser-place}

Three common laser placements:

\begin{description}
\item[On-engine / on-PIC] Shortest optical path; hottest and least
      serviceable. Attractive when reliability and thermal headroom are
      proven.
\item[On-package but separated] Laser die near the engine with some
      thermal isolation. Compromise on coupling length versus heat.
\item[External (ELS / ELSFP)] Lasers at the cool faceplate, fiber or
      waveguide into the package. Best serviceability and laser ambient;
      adds coupling, connectors, and shared-source failure domains
      (\Cref{sec:elsfp}).
\end{description}

Interview line: ``I choose laser placement from thermal life and field
replaceability first, then from coupling loss. Power per bit alone does not
pick the architecture.''

\section{Placement ladder}
\label{sec:pkg-placement-ladder}

Read the optics-placement ladder once as a design menu, not as a chronology
you must climb:

\begin{center}
\footnotesize
\begin{tabular}{@{}ll@{}}
\toprule
Step & Design question \\
\midrule
Pluggable & Can faceplate serviceability and VSR SI still close? \\
OBO / COBO & Is mid-board optics worth losing hot-swap? \\
NPO & Engine beside ASIC on substrate: enough reach cut? \\
CPO & XSR die-to-engine worth package yield and thermal risk? \\
Interposer optical I/O & On-package photonics with almost no copper left? \\
\bottomrule
\end{tabular}
\end{center}

\Cref{tab:placement} carries energy-per-bit and serviceability numbers.
\Cref{sec:cpo-status} carries who is shipping what. This chapter asks whether
your constraints match the step you named.

Copper counter-moves (CPC / NPC) push the crossover outward for one more rate
generation (\Cref{app:ai-fabric-context}). They do not erase the packaging
trade when reach or port count forces optics inward.

\section{Thermal and power design for optical engines}
\label{sec:pkg-thermal}

Package thermal design is not a heatsink footnote.

\begin{itemize}
\item Map ASIC hotspot, optical-engine location, and laser location on one
      sketch.
\item Budget junction temperature for lasers separately from case
      temperature for the module (\Cref{ch:lasers}).
\item Treat microring heaters as both actuators and heat sources
      (\Cref{ch:wdm}).
\item Liquid cooling at the rack does not automatically cool a laser buried
      under a lid next to a switch die.
\item Power delivery noise and simultaneous switching can look like optical
      BER if supplies are shared carelessly.
\end{itemize}

\Cref{sec:thermal-envelope,sec:power} expand fabric-level envelopes. Here the
rule is local: if the optical engine sees the ASIC thermal gradient, wavelength
and life margins shrink even when the room looks fine.

\section{Known-good die, yield, and test before assembly}
\label{sec:pkg-kgd}

A whiteboard packaging answer that skips test is incomplete.

\begin{enumerate}
\item Screen ASICs and optical engines as far as probe, known-good-die, and
      partial assembly allow.
\item Define optical attach metrics before lid seal: coupling loss, dark
      lanes, polarization, monitor PD continuity.
\item Keep a rework path where the business case allows it. Pure 3D stacks
      often do not.
\item Plan system-level bring-up that can still localize package versus host
      versus fiber plant (\Cref{ch:productization,ch:failure-modes}).
\end{enumerate}

Package yield is multiplicative when you stack. Separate screening is why many
CPO engines stay 2.5D-adjacent rather than fully bonded until the optical test
story matures.

\section{Whiteboard vignette}
\label{sec:pkg-vignette}

Prompt: ``Design the packaging for a 102.4~Tb/s switch optical shoreline.
Do we go CPO? 2.5D or 3D? Where do the lasers live?''

A strong answer stays ordered:

\begin{enumerate}
\item Name the binding constraint: SerDes energy, faceplate density, or
      serviceability.
\item If CPO is required, pick 2.5D-adjacent engines unless a bonded 3D
      electronics-photonics stack has a proven KGD and thermal path.
\item Put lasers in ELSFP-class banks unless on-package laser life is already
      proven at the ASIC hotspot.
\item Define FAU attach screens and dark-lane reject before lid seal.
\item Co-design XSR, power rails, and heatsink keep-outs with the ASIC team.
\item State the field replaceables explicitly for the reliability model.
\end{enumerate}

Weak answers jump to ``use COUPE 3D'' without naming yield and service. Strong
answers treat vendor process names as options inside a constraint set.

\section{Interview takeaway}
\label{sec:pkg-design-review}

\keyidea{I name the placement first: pluggable, NPO, or CPO. Then I name the
integration style: 2.5D or 3D. I do not treat those as synonyms. I partition
laser heat from the ASIC, insist on a known-good-die and fiber-attach story,
and I only accept soldered engines when field replacement of the laser bank
or a clear non-serviceable cost model closes.}

\section{Interview Q\&A}
\label{sec:pkg-interview-qa}

Practice aloud. Prefer first-person reasoning. Score with
\Cref{sec:chapter-interview-rubric}.

\paragraph{Question 1. What is the difference between 2.5D and 3D?}
\textit{Tests:} lateral versus vertical integration.

\textit{Spoken answer.}
``2.5D puts dies side by side on an interposer, bridge, or RDL, so the dense
wiring is lateral. 3D stacks dies with hybrid bonds, microbumps, or TSVs, so
the dense wiring is vertical. 3D wins bandwidth density and pays in thermal
access, test access, and stack yield.''

\textit{Pressure follow-up.} ``Is CPO the same as 3D?''\\
\textit{Answer pivot.} ``No. CPO is where the optical engine sits relative to
the ASIC. It can be 2.5D, 3D, or substrate-mounted.''

\textit{Trap:} ``CPO means hybrid-bonded 3D photonics.''

\paragraph{Question 2. Why choose CPO?}
\textit{Tests:} electrical reach and energy density.

\textit{Spoken answer.}
``I choose CPO when faceplate VSR copper and pluggable power stop closing the
port-count or energy budget. Shortening to XSR cuts SerDes energy and board SI
pain, and it raises shoreline density. I still have to close thermal, yield,
fiber, and serviceability.''

\textit{Pressure follow-up.} ``Why not just use LPO pluggables?''\\
\textit{Answer pivot.} ``LPO deletes module DSP but keeps the long host
electrical run and cage. CPO attacks the copper reach itself.''

\textit{Trap:} ``CPO is always lower power, so always choose it.''

\paragraph{Question 3. Why might you refuse CPO?}
\textit{Tests:} yield, thermal, serviceability judgment.

\textit{Spoken answer.}
``I refuse CPO when package yield is unproven, when the laser cannot be kept
off the ASIC hotspot, when fiber attach scrap risk is unacceptable, or when
the fleet requires faceplate hot-swap of the whole optical path. XPO or
pluggables can be the rational hold.''

\textit{Pressure follow-up.} ``Management says competitors shipped CPO.''\\
\textit{Answer pivot.} ``Shipping proves a recipe can work. It does not prove
our thermal partition, FAU yield, or service model are ready.''

\textit{Trap:} ``Never do CPO until every vendor program looks perfect.''

\paragraph{Question 4. External versus integrated laser?}
\textit{Tests:} thermal life versus coupling and failure domain.

\textit{Spoken answer.}
``An integrated laser shortens the optical path but sits near the heat and is
hard to replace. An ELSFP-class external laser keeps the source cool and
field-replaceable, at the cost of coupling, connectors, and a shared-source
failure domain. I pick from life and service first, then from coupling loss
(\Cref{sec:elsfp}).''

\textit{Pressure follow-up.} ``Does external laser make CPO unnecessary?''\\
\textit{Answer pivot.} ``No. It often enables CPO by removing the least
reliable part from the hot package.''

\textit{Trap:} ``External lasers are only for discrete pluggables.''

\paragraph{Question 5. How do you partition thermally?}
\textit{Tests:} hotspot map and wavelength/life coupling.

\textit{Spoken answer.}
``I sketch ASIC hotspot, engine, laser, and fiber exit. I keep lasers off the
hottest region when life matters, treat ring heaters as heat sources, and I
budget junction temperature separately from case temperature. Rack liquid
cooling does not automatically fix an on-die thermal gradient.''

\textit{Pressure follow-up.} ``The case temperature meets spec. Why are rings
walking?''\\
\textit{Answer pivot.} ``Local gradient and neighbor heaters can move
resonances even when case average looks fine.''

\textit{Trap:} ``If the heatsink is sized for ASIC watts, optics are covered.''

\paragraph{Question 6. What is known-good die in this context?}
\textit{Tests:} screen before irreversible attach.

\textit{Spoken answer.}
``Known-good die means I screen the ASIC and the optical engine as far as
probe and partial assembly allow before I commit to a stack or sealed
co-package. A bad PIC bonded to a good switch wastes both. 3D makes this
harder because stack yield multiplies.''

\textit{Pressure follow-up.} ``Can you fully optically test at wafer probe?''\\
\textit{Answer pivot.} ``Often not completely. I define which optical metrics
are probeable, which need partial assembly, and which wait for FAU attach.''

\textit{Trap:} ``Wafer sort of the ASIC is enough; optics will be caught in
system test.''

\paragraph{Question 7. How does package yield change the architecture?}
\textit{Tests:} multiplicative yield and cost of scrap.

\textit{Spoken answer.}
``If optical attach or fiber alignment yield is soft, I avoid architectures
that scrap a whole switch package on every miss. That pushes toward separable
engines, reworkable attach, or keeping lasers in replaceable modules. Yield
is an architecture input, not only a factory metric.''

\textit{Pressure follow-up.} ``Yield is a manufacturing problem.''\\
\textit{Answer pivot.} ``Manufacturing executes it. Architecture chooses
whether a fail scraps \$X or \$100X.''

\textit{Trap:} ``Ship the densest stack and let the factory climb the yield
curve.''

\paragraph{Question 8. What matters for optical I/O in the package?}
\textit{Tests:} coupling style and keep-outs.

\textit{Spoken answer.}
``I name edge versus grating coupling, fiber exit direction, polarization
control, and lid or socket keep-outs. Optical I/O has to coexist with the
heatsink, power planes, and board clearance. A beautiful electrical 3D stack
that cannot exit fiber cleanly is not a product.''

\textit{Pressure follow-up.} ``Can we always use grating couplers for easier
vertical fiber?''\\
\textit{Answer pivot.} ``Sometimes. They trade efficiency, polarization, and
wavelength dependence. Edge coupling may win on loss if the mechanical design
allows it.''

\textit{Trap:} ``Optical I/O is just a FAU drawing detail after electrical
floorplanning.''

\paragraph{Question 9. Why is fiber attach a first-class risk?}
\textit{Tests:} coupling, dark lanes, sealed-package discovery.

\textit{Spoken answer.}
``FAU attach sets insertion loss, lane balance, and often polarization. A
single misaligned channel looks like a dead laser or bad PIC after the lid is
on. I want attach metrics and dark-lane screens before irreversible seal, and
I track FAU variation in FA when one lane is weak
(\Cref{sec:photonic-packaging}).''

\textit{Pressure follow-up.} ``Average coupling looks fine.''\\
\textit{Answer pivot.} ``I still check the weakest lane. Averages hide the
channel that fails the link.''

\textit{Trap:} ``If total power out of the FAU is in family, attach is good.''

\paragraph{Question 10. How do you think about repairability?}
\textit{Tests:} field replaceable elements versus soldered engines.

\textit{Spoken answer.}
``I list what can fail and what a tech can replace: ELSFP laser bank, fiber
jumper, pluggable module, or nothing because the engine is soldered. CPO often
keeps the engine soldered and makes the laser replaceable. I state that
explicitly so reliability and ops models match the hardware.''

\textit{Pressure follow-up.} ``Is a soldered engine always a bad idea?''\\
\textit{Answer pivot.} ``No, if the dominant fails are elsewhere and the cost
model accepts package replacement. It is a bad idea if laser FIT still
dominates and lasers are trapped on the hot die.''

\textit{Trap:} ``CPO cannot be serviced, so it cannot ship.''

\paragraph{Question 11. What does switch-ASIC co-design mean here?}
\textit{Tests:} shoreline, SerDes, power, and thermal co-ownership.

\textit{Spoken answer.}
``The optical engine and the switch SerDes share shoreline, bump map, power
delivery, and thermal path. I co-design lane count, XSR budget, retimer or
DSP placement, and heatsinking. A perfect PIC that the ASIC cannot drive
cleanly, or that cooks the rings, is not an engine.''

\textit{Pressure follow-up.} ``Who owns the XSR budget?''\\
\textit{Answer pivot.} ``Both sides. Package SI and SerDes EQ share it
(\Cref{ch:equalization,sec:serdes-dsp}).''

\textit{Trap:} ``The optics team owns optics; the ASIC team owns silicon;
package is NPI's problem.''

\paragraph{Question 12. How do you test an optical engine before final
assembly?}
\textit{Tests:} staged test plan.

\textit{Spoken answer.}
``I stage tests: die or engine probe where possible, electrical continuity and
basic photonic metrics after attach, optical power and dark-lane screens after
FAU, then full link metrics after lid and at system bring-up. Each stage needs
pass/fail that prevents carrying a known-bad unit into a more expensive
assembly step.''

\textit{Pressure follow-up.} ``System BER is the only test that matters.''\\
\textit{Answer pivot.} ``System BER is necessary, but it is a late, expensive
screen. Earlier stages protect yield.''

\textit{Trap:} ``Build the full package, then characterize everything at the
host.''

\paragraph{Question 13. Give a 60-second packaging plan for a CPO engine.}
\textit{Tests:} end-to-end judgment order.

\textit{Spoken answer.}
``I start from reach and port density to decide if CPO is needed. I choose
2.5D versus 3D from thermal access, KGD, and optical I/O, not from buzzwords.
I place lasers for life and service, often external. I define FAU attach and
dark-lane screens before seal, co-design XSR and power with the ASIC, and I
only freeze the package when yield, thermal corners, and field replacement
match the fleet model.''

\textit{Pressure follow-up.} ``Schedule is cut. What do you protect?''\\
\textit{Answer pivot.} ``Laser thermal partition, FAU yield evidence, and a
service story for the highest-FIT part. Density slides after those.''

\textit{Trap:} ``Pick the densest COUPE-like stack and qualify later.''
